% !PRINTED 623
\weberpage{623}
Bilden wir aus (2) die $p^{\text{te}}$ Potenz von $\omega$, und beachten den
\gesperrt{Fermat'schen} Lehrsatz f\"ur rationale Zahlen $[a_0^p \equiv a_0 \pmod{p}]$,
so ergiebt sich
\[
\omega^p \equiv a_0 + a_1 r^p + a_2 r^{2p} + \cdots + a_{\mu-1} r^{(\mu-1)p} \pmod{p},
\]
oder nach (2)
\weq{3}{\omega^p \equiv \omega_p \pmod{p}.}

Dadurch ist die Gruppe $\Psi$ bestimmt, zu der das Ideal $\mathfrak{p}$
geh\"ort. Es ist n\"amlich nach (3) auch
\[
\omega^p \equiv \omega_p \pmod{\mathfrak{p}},
\]
daher ist $\psi_0 = (r, r^p)$ und die Gruppe $\Psi$ besteht nach \S.\,162, (13)
aus den Potenzen dieser Substitution, soweit sie von einander
verschieden sind.

Ist daher $f$ der kleinste positive Exponent, f\"ur den
\weq{4}{p^f \equiv 1 \pmod{m}}
ist, d.\,h. geh\"ort $p$ zu dem Exponenten $f$ f\"ur den Modul $m$, so ist
\weq{5}{\Psi = (r,r), (r,r^p), (r,r^{p^2}), \ldots (r,r^{p^{f-1}})}
und ist vom $f^{\text{ten}}$ Grade. Demnach ist auch $\mathfrak{p}$ ein Primideal
$f^{\text{ten}}$ Grades und
\weq{6}{N(\mathfrak{p}) = p^f.}

Die Anzahl $e$ der von einander verschiedenen conjugirten
Primfactoren von $f$ ist $\mu : f$, und wir haben also den Satz: % [sic] printed $f$ after "von"; mathematically $p$ is required

III. \gesperrt{Ist $p$ eine von $q$ verschiedene Primzahl, die f\"ur den Modul $m$ zum Exponenten $f$ geh\"ort, ist ferner $\mu = \varphi(m) = ef$, so zerf\"allt $p$ in $\Omega_m$ in $e$ von einander verschiedene conjugirte Primfactoren $f^{\text{ten}}$ Grades.}

\begin{center}
\S.\,172.

\gesperrt{Die conjugirten Primideale.}
\end{center}

Wir m\"ussen noch etwas genauer auf die Bildung der conjugirten Primfactoren $\mathfrak{p}_n$ einer von $q$ verschiedenen Primzahl $p$
eingehen, die nach dem zuletzt bewiesenen Satze in $\Omega_m$ in $e$
von einander verschiedene Primfactoren zerf\"allt, wenn
\weq{1}{\varphi(m) = ef}
ist, und $f$ der kleinste positive Exponent, f\"ur den
\weq{2}{p^f \equiv 1 \pmod{m}.}
